\section{Toward a Stable Multi-Task OPD Recipe}
\label{sec:recipe}

The three studies suggest that capability retention should guide the choice of averaging rule, teacher feedback, and training checkpoint. A useful starting point is to specify the acceptable tradeoff before training. Let $s_d(\theta)$ be a development-set score for domain $d$, oriented so that higher is better, and let $\epsilon_d\geq0$ be a prespecified tolerance relative to the initial student. The proposed recipe selects checkpoints subject to
\begin{equation}
s_d(\theta)\geq s_d(\theta_0)-\epsilon_d
\qquad\text{for every domain }d.
\label{eq:retention}
\end{equation}
Among checkpoints meeting these conditions, the proposed rule selects the one with the highest prespecified aggregate utility, retaining the initial student if none improves it. Selection should account for score uncertainty at the chosen tolerances, followed by evaluation on a separate test set. The constraints govern checkpoint selection rather than add differentiable terms to the distillation loss.

With the retention goal fixed, the averaging rule should express the intended importance of each training response. For an objective that weights responses equally within each domain, we use DR with explicit domain weights $\lambda_d$. In our four-domain setting, equal weights and 16 prompts per domain provide the starting allocation. If longer responses are deliberately assigned more weight within a domain, DT expresses that alternative. Under either reduction, the prompt quota and the domain loss weight remain separate choices: with every domain represented and its loss averaged separately, increasing a quota at fixed $\lambda_d$ changes exposure and estimation variance without directly increasing that domain's expected gradient weight. Generated lengths and truncation rates should therefore be recorded alongside both quantities when diagnosing an imbalance.

Once these weights are explicit, vocabulary feedback can be chosen using its capability profile and computational cost. The joint comparison in Section~\ref{sec:density} motivates I64 as a candidate alongside DR, with PG retained as a reference. Its feedback scale also deserves attention. For the student candidate set $S$ and teacher intersection $I$ in Equation~\ref{eq:i64}, the retained student weights sum to $r(h)=p_\theta(I\mid h)/p_\theta(S\mid h)$. At a fixed prefix with a nonempty intersection and detached coefficients, the I64 gradient is $r(h)$ times the gradient obtained by renormalizing the same weights over $I$. Thus removing length weighting does not remove variation in feedback scale across prefixes. Monitoring this retained mass and empty intersections by domain helps assess whether a fixed candidate size supplies comparable coverage; increasing or renormalizing a small intersection requires validation because it can also amplify unreliable feedback.

If a domain repeatedly falls below its retention tolerance, the response should depend on the diagnosed cause. Exposure, length, and coverage checks can reveal inadequate sampling or unintended weighting. When these do not explain the decline, teacher interactions should be examined at a common student and optimizer state. The overlap diagnostic defined in Section~\ref{sec:subnetworks} describes shared update locations; a further functional check is to apply each teacher's proposed step and measure the change in fixed probe losses from all domains. A zero-current-gradient step provides a reference for the contribution of optimizer history. This comparison can identify local cross-domain harm that coordinate overlap alone does not establish, while subsequent capability evaluations determine whether it persists along training. Gradient modification or task-specific parameters become interventions to test when such harm remains reproducible after the allocation and coverage checks.

The default recipe keeps domain weights and prompt quotas fixed; any adaptive variant should state its trigger and adjustment rule separately. To maintain previous-domain exposure without changing the on-policy protocol, retain prompts from those domains and generate fresh responses with the current student before teacher scoring. The resulting starting recipe combines explicit response-level weighting, I64 feedback with coverage monitoring, and retention-based checkpoint selection. Its stability must be assessed both along the training trajectory and across independent runs: an acceptable endpoint does not establish that earlier checkpoints avoided degradation.

\completionnote{Recipe validation}{Complete a matched GT/DR $\times$ PG/I64 comparison through update 250 with at least three paired training seeds, adding PG--GT to the existing configurations. Retain I64--DT to separate between-domain and within-domain length effects. Match per-domain prompt exposure and update counts, report token and compute costs separately, and evaluate domain trajectories, worst-domain changes, and retention across seeds. Use an independent development set for checkpoint selection. Complete the common-state teacher-step comparisons with cross-domain probe losses and a zero-current-gradient reference, and audit I64 retained mass by domain. These comparisons are pending; no aggregate recipe result is reported here.}
